% Introduction
% Structure: Problem/Motivation -> Background -> Research Questions -> Contributions -> Paper Outline

% Motivating example or opening
Wikipedia, one of the world’s largest collaborative knowledge platforms, operates under an open-editing model that invites anyone to contribute. While this openness has democratized knowledge creation, it also raises complex questions about bias, neutrality, and governance. For example, discussions on politically charged articles such as \textit{Climate Change} or \textit{Gun Politics in the United States} often reveal tensions between consensus-building and ideological polarization among editors.

% The problem/question
The central problem this research addresses is how Wikipedia’s open, collaborative structure influences the neutrality of its content and the democratic quality of its editorial deliberations. Specifically, we examine whether Wikipedia’s design fosters genuinely democratic decision-making—or whether, as prior research suggests, governance tends to consolidate among small groups of experienced editors.

% Why it matters
Understanding how Wikipedia handles bias and governance is critical not only for evaluating the reliability of one of the most visited knowledge resources online, but also for informing the design of future collaborative systems. As Wikipedia continues to shape global discourse, identifying how bias emerges, evolves, and is moderated through collective deliberation offers broader insights into the health of digital democracy.

% Research questions
This paper investigates the following research questions:
\begin{enumerate}
    \item How can we detect and measure biased language in Wikipedia articles at the sentence level?
    \item How does Wikipedia bias change over time through the revision process?
    \item To what extent do Wikipedia Talk pages function as democratic spaces of deliberation?
    \item How do patterns of participation influence consensus-building?
    \item Does democratic participation reduce bias, or do small groups of dominant editors shape outcomes that increase bias?
\end{enumerate}

% Contributions
The main contributions of this work are:
\begin{itemize}
    \item A computational framework for detecting and tracking bias in Wikipedia articles using sentence-level machine learning models.
    \item An empirical analysis of Talk page discussions to assess sentiment, toxicity, and deliberative tone across topics.
    \item An integrative model connecting editorial diversity and governance patterns with the neutrality of Wikipedia content.
    \item A prototype system leveraging the Wikipedia API for automated analysis and visualization of editor interactions.
\end{itemize}

% Paper outline
The rest of this paper is organized as follows. 
Section~\ref{sec:related} reviews related work on Wikipedia governance, bias, and online deliberation. 
Section~\ref{sec:methodology} describes our data sources and analytical methods, including bias detection and governance modeling. 
Section~\ref{sec:results} presents findings from the prototype implementation and sentiment/toxicity analyses.
Section~\ref{sec:discussion} interprets the implications of these results for digital democracy and collaborative systems.
Section~\ref{sec:conclusion} concludes and outlines directions for future work.
